\scanpage{101}
qui pour $S$ variable est compatible aux chgts de base, et donne donc un morphisme des multiplicités modulaires \inserted{(si $G,G'$ sont stables)}
\[
(132)\qquad M_G\longrightarrow M_{G'},
\]
avec fonctorialité évidente pour les composés. Ainsi $M_G$ est un \uline{foncteur contravariant en} $G$ pour les morphismes de normalisation \inserted{de m.graphes stables}. Cette fois, on a

\uline{Proposition} a) Si $G\longrightarrow G'$ est un morphisme de normalisation \uline{strict}, alors (132) est un \uline{iso.} des multiplicités

b) Soient $G_\alpha$ des \deleted{\unsure{graphes connexes}}\inserted{m.graphes stables} \deleted{de} $G$ le m.graphe somme des $G_\alpha$. Alors le morphisme
\[
(133)\qquad M_G\longrightarrow\prod_\alpha M_{G_\alpha}
\]
est un iso. des multiplicités.
\margin{\uline{NB} On remplace la notation $\Pi M_G$ par $\Pi_1M_G$, pour ne pas confondre $\unsure{\Pi|\Pi_1}\deleted{\unsure{\Pi}}$ avec le signe du produit.}

\uline{Cor} $\Pi_1M_G$ est un foncteur contravariant en \deleted{\unsure{le}} m.graphe stable $G$ \deleted{\unsure{f}} pour les
\margin{100}
